\section{Circuit Design}
\paragraph{Circuit Types}
We focus on digital circuits design with the advantages over analog circuits:
\begin{itemize}
  \item More reliable, with simpler circuits and less noise-prone
  \item Specified accuracy (determinable)
  \item Abstraction can be applied using boolean algebra
  \item Simpler design, analysis and simplification
\end{itemize}
Among digital circuits, we have \textit{combinational circuits} (with no
memory, and output depends solely on the input) and \textit{sequential
circuits} (where output depends on both input and current state of memory).

\subsection{Boolean Algebra}
\subsubsection{Operations}
\begin{notation}
  In Boolean Algebra, True is denoted as $T$ or $1$, and False as $F$ or $0$.
  We denote conjunction (AND) as $A\cdot B$ or $A\wedge B$, disjunction (OR) as
  $A+B$ or $A\vee B$, and negation (NOT) as $A'$, $\bar{A}$ or $\neg{A}$. 
  \begin{remark}
    In CS2100, we use $1$ for True, $0$ for false, $\cdot$ for conjunction, $+$
    for disjunction and $'$ for negation.
  \end{remark}
\end{notation}
\paragraph{Precedence of Operators}
There is no convention for boolean algebra in general. We define the precedence
of operators from higher to lowest (for CS2100):
\begin{enumerate}
  \item Not ($'$)
  \item And ($\cdot$)
  \item Or ($+$)
\end{enumerate}
\subsubsection{Results of Boolean Algebra}
We have the following laws of boolean algebra:

\centering
\begin{tabular}{c c c}
  \toprule
  Law & Disjunction & Conjunction\\
  \midrule
  Identity Laws & $A+0=0+A=A$ & $A\cdot 1=1\cdot A = A$ \\
  Inverse/Complement Laws & $A+A'=A'+A=1$ & $A\cdot A'=A'\cdot A=0$ \\
  Commutative Laws & $A+B=B+A$ & $A\cdot B=B\cdot A$ \\
  Associative Laws & $A+(B+C)=(A+B)+C$ & $A\cdot(B\cdot C)=(A\cdot B)\cdot C$ \\
  Distributive Laws & $A\cdot (B+C)=(A\cdot B)+(A\cdot C)$ & $A+(B\cdot C)=(A+B)\cdot(A+C)$ \\
  \bottomrule
\end{tabular}\\
\vspace{4mm}
\raggedright%
and the following theorems:\\
\vspace{2mm}

\centering
\begin{tabular}{c c c}
  \toprule
  Theorem & Disjunction & Conjunction \\
  \midrule
  Idempotency & $X+X=X$ & $X\cdot X=X$ \\
  One Element / Zero Element & $X+1=1$ & $X\cdot 0 = 0$ \\
  Involution & \multicolumn{2}{c}{$\left( X' \right)'=X $ } \\
  Absorption 1 & $X+X\cdot Y=X$ & $X\cdot (X+Y)=X$ \\
  Absorption 2 & $X+X'\cdot Y=X+Y$ & $X\cdot (X'+Y)=X\cdot Y$ \\
  De Morgan's Laws & $(X+Y)'=X'\cdot Y'$ & $(X\cdot Y)'=X'+Y'$ \\
  \multirow{2}{4em}{Consensus} & \multicolumn{2}{c}{$X\cdot Y+X'\cdot Z + Y\cdot Z=X\cdot Y+X'\cdot Z$}\\
  & \multicolumn{2}{c}{$(X+Y)\cdot(X'+Z)\cdot(Y+Z)=(X+Y)\cdot(X'+Z)$} \\
  \bottomrule
\end{tabular}

\raggedright%
\paragraph{Duality} Each result comes in pairs. If each AND/OR operator and
identity element 0/1 in a boolean equation is interchanged, it remains valid. This
is known as \textit{duality}.

\subsubsection{Standard Forms}
\paragraph{Definitions}
\begin{definition}[Terms]
  A boolean variable, or its complemented form is known as a \textit{literal}.
  A single literal, or a logical product (AND) of literals is a \textit{product
  term}. Similarly, a single literal, or a logical sum (OR) of literals is a
  \textit{sum term}.
\end{definition}
We define two standard forms of boolean expressions, to the end of desirable
circuits for implementation:
\begin{itemize}
  \item \textit{Sum-of-Products (SOP)}: A product term, or a logical sum (OR) of several product terms.
  \item \textit{Product-of-Sums (POS)}: A sum term, or a logical product (AND) of several sum terms.
\end{itemize}
\begin{remark}
  All boolean expressions can be expressed in SOP and POS form.
\end{remark}
\begin{definition}[Minterm, Maxterm]
  A \textit{minterm} of $x$ variables is a product term containing $x$
  literals, denoted m0 to m$(2^{n}-1)$. A \textit{maxterm} of $x$ variables is a
  sum term containing $x$ literals, denoted M0 to M$(2^{n}-1)$.
\end{definition}
\begin{example}[Number of min-, maxterms]
  Given $n$ variables, we have up to $2^{n}$ minterms and $2^{n}$ maxterms.
  \begin{figure}[h!]
    \centering
    \includegraphics[width=0.5\textwidth]{minMaxTerms.png}
  \end{figure}
\end{example}
\textbf{Each minterm is the complement of the corresponding maxterm.}
To find the m(x) minterm, where $x={(b_1b_2\cdots b_{n})}_2$, we have the
product of $n$ literals of which each literal $a_i=\begin{cases}a_i'&b_i=0\\a_i&b_i=1\end{cases}$. The M(x) maxterm is then the sum of $n$ literals where
each literal $a_i =\begin{cases}a_i&b_i=0\\a_i'&b=1\end{cases}$

Using minterms and maxterms, we may then obtain a \textbf{canonical/normal
form}, which ensures uniqueness of the SOP or POS\@.

A \textit{sum-of-minterms} is the canonical SOP, denoted $\sum m(k_1,k_2,\ldots
k_i)$. It can be obtained by gathering the minterms with which the output is 1.
A \textit{product-of-maxterms} is the canonical POS, denoted $\prod
M(k_1,k_2,\ldots,k_i)$. It can be obtained by gathering the maxterms with which
the output is 0.

Additionally, each function $F=\sum m(k_1,k_2,\ldots,k_i)=\prod M(l_1,l_2,\ldots,l_j)$
where $\left\{ 0,\ldots,n-1 \right\}\setminus\left\{ k_1,k_2,\ldots k_i
\right\}=\left\{ l_1,l_2,\ldots l_j \right\} $.


\subsection{Circuit Simplification}
Given a boolean expression, we may then simplify to obtain a circuit that uses
fewer logic gates, such that it is cheaper, uses less power and (sometimes)
faster. We do this with various simplification techniques, namely
\textbf{algebraic simplification}, \textbf{Karnaugh Maps}, and
the \textbf{Quine-McCluskey Technique}.

\subsubsection{Algebraic Simplification}
By application of the boolean algebra results and theorems, we are able to
minimise the number of literals and number of terms. However, these two
minimisations may sometimes be in conflict, in which case we have to opt to
adapt the tradeoff within the constraints in our context. \\

Algebraic manipulation allows us to adapt the technique, but is sometimes
challenging, requiring strong algebraic manipulation skills.

\subsubsection{Karnaugh Maps}
\begin{definition}[Gray Code]
  An ordering of the binary numerals is a \textit{Gray code} if every two
  successive values differ in only one bit.
\end{definition}
A Gray code is an unweighted code, and every of the $n$-bits are usually
utilised to have a sequence of $2^n$ values. The $n$-bit standard gray code can
be obtained inductively from the $(n-1)$-bit standard gray code.
\begin{figure}[h!]
  \centering
  \includegraphics[width=0.26\textwidth]{gray_code}
\end{figure}\\
Given $n$-bits, there are multiple possible Gray code cycles, from which the
standard Gray code uniquely specifies one. For instance, we have the $4$-bit
standard gray code as shown.
\begin{figure}[h!]
  \centering
  \includegraphics[width=0.5\textwidth]{4_gray_code}
\end{figure}\\
Within a universal set $U$ of minterms, we can partition the minterms, from
which a boolean function is a subset where we obtain the SOP expression. We may
use venn diagrams to analyse boolean expressions.
\begin{definition}[Karnaugh Maps]
  A \textit{K-Map} is an abstract form of a Venn diagram organised as a matrix
  of cells where
  \begin{itemize}
    \item Each cell represents a minterm
    \item Adjacents cell differ by one literal
  \end{itemize}
\end{definition}
We define K-maps up to 6 variables. For each K-map, we fill each cell with $1$
if it corresponds to a minterm of the function, $d$ if it corresponds to a
dont-care term, and $0$ otherwise. Given an arbitrary Boolean function, we can
convert it to a SOP form, where we can then fill the K-map directly, or
construct the truth table to obtain the minterms.
\paragraph{n-Variable K-Maps} Given n variables $a,b$, we have the K-maps as
shown below, with axes each Gray codes representing the status of the
variables, and where each cell has n neighbours.
\begin{figure}[h!]
  \centering
  \includegraphics[width=0.15\textwidth]{kmap2}
  \includegraphics[width=0.25\textwidth]{kmap3}
\end{figure}
\begin{figure}[h!]
  \centering
  \includegraphics[width=0.4\textwidth]{kmap4}
\end{figure}\\
\begin{figure}[h!]
  \centering
  \includegraphics[width=0.73\textwidth]{kmap5}
\end{figure}
\begin{figure}[h!]
  \centering
  \includegraphics[width=0.5\textwidth]{kmap6}
\end{figure}
\vspace{16mm}
Given a filled K-map, we may obtain valid groupings of adjacent cells as follows:
\begin{itemize}
  \item Each cell may be grouped with any of its neighbouring cell
  \item Each group has size in powers of two
  \item Each group correspond to a product term
  \item Each don't-care term can be taken to be $1$ or $0 $ to obtain a simpler
    expression
\end{itemize}
\begin{remark}
  Each grouping of $2^{n}$ cells eliminates $n$ variables from the product
  terms.
\end{remark}
\begin{definition}[Prime Implicants]
  A product term that is obtained by combining the maximum number of minterms
  from adjacent squares in the map is known as a \textit{prime implicant}.
\end{definition}
\begin{definition}[Essential Prime Implicant]
  A prime implicant that includes at least one minterm that is not covered by
  any other prime implicant is known as a \textit{essential prime implicant}.
\end{definition}
Given the K-map, we may then find a \textbf{simplified SOP expression} to the
function with the following algorithm.
\begin{enumerate}
  \item Identify PIs of the function by finding all largest possible groupings.
  \item Within the PIs, identify the EPIs that must exist in the expression.
  \item Select a minimum subset of the remaining PIs to complete the cover.
\end{enumerate}

Alternatively, we may find a \textbf{simplified POS expression} to the function
by grouping the maxterms. This is equivalent to grouping the minterms in the
K-map of the function's complement. We can then obtain the POS expression by
taking the inverse of the SOP expression found from this grouping.\\

By using K-Maps, we can minimise the number of product terms and literals. They
are easy to use, but limits us to 5 to 6 variables.


\subsubsection{Quine-McCluskey Technique}
With $n$ variables, in the Quine-McCluskey Technique, we denote each product
term as a $n$-length string, with $1, 0$ for the literals, and $-$ for an
absence of a variable. We can obtain a simplified expression using an iterative
algorithm. Given a sum-of-minterms expression, 
\begin{enumerate}
  \item Arrange the terms in groups according to the number of $1$s.
  \item Combine terms from adjacent groups which differ by only $1$ bit.
  \item Iterate step $2$ until no further combinations are possible.
\end{enumerate}
When terms are combined, we indicate them as such with a tick. When the program
terminates, terms that have not been combined are the PIs. We then find a cover
using a PI chart, to find the EPIs given by columns with a single $X$, and then
a reduced PI chart to find a minimum cover for the remaining minterms.
\begin{example}
  Given a $4$-variable function
  \[F(A,B,C,D)=\sum m\left( 2,4,6,8,9,10,12,13,15 \right)\] 
  we have the following table from applying the algorithm.
  \begin{figure}[h!]
    \centering
    \includegraphics[width=0.5\textwidth]{quine_mccluskey}
  \end{figure}\\
  From this, we have the PI chart and then the reduced PI chart.
  \begin{figure}[h!]
    \centering
    \includegraphics[width=0.4\textwidth]{quine_pi}
    \includegraphics[width=0.3\textwidth]{quine_reducedpi}
  \end{figure}
  As such, we have the minimal SOP expression
  \[F=A\cdot C'+A\cdot B\cdot D+B'\cdot C\cdot D+A'\cdot B\cdot D'.\] 
\end{example}
This technique guarantees to produce a minimal standard form, and can be
applied to functions on many variables. However, it is highly tedious and more
suited for automation.


\subsection{Logic Circuits}
A logic circuit is composed of logic gates, represented as shown. We call the number of inputs of a gate the \textbf{fan-in}. Gates can have a fan-in more than 2.
\begin{figure}[h!]
  \centering
  \includegraphics[width=0.5\textwidth]{logicGates.png}
\end{figure}
\begin{remark}
  By De Morgan's laws, we have the NAND $\equiv$ Negative-OR gate, and NOR $\equiv$ Negative-AND.
\end{remark}
Given a boolean expression, we may implement it as a logic circuit.
Additionally, given a logic circuit, we may obtain its logic expression.
We call the set $\left\{ AND, OR, NOT \right\} $ a \textbf{complete set of logic}. It is
sufficient for building any boolean function.
\begin{definition}[Universal Gates]
  A \textbf{Universal Gate} is a gate that itself is a complete set of logic. We may
  prove a specified gate is a universal gate by implementing a complete set of
  logic with the gate.
\end{definition}
\begin{example}[NAND]
  NAND is a universal gate.
  \begin{figure}[h!]
    \centering
    \includegraphics[width=0.25\textwidth]{nandUniversal}
  \end{figure}
\end{example}
\begin{example}[NOR]
  NOR is a universal gate.
  \begin{figure}[h!]
    \centering
    \includegraphics[width=0.25\textwidth]{norUniversal}
  \end{figure}
\end{example}
We may implement a given a SOP expression using a 2-level AND-OR circuit or a
2-level NAND circuit.
\begin{proof}
  With a SOP expression $F=A_0\cdot A_1 \cdot\ldots \cdot A_{n}+\cdots+Z_0\cdot
  Z_1\cdot\ldots\cdot Z_{n}$, we obtain each product term with AND gates for
  each product term, and obtain the final expression with a single OR gate.  By
  negating the output of each AND gate, and then negating the input to the OR
  gate, we obtain a 2-level circuit of NAND gates.\qed
\end{proof}\\
Similarly, we may implement a given a POS expression using a 2-level OR-AND
circuit or a 2-level NOR circuit.

\subsubsection{Propagation Delay}
\begin{definition}[Propagation Delay]
  The \textbf{gate delay} or \textbf{propagation delay} of a given gate is the
  average transition time taken for the output signal of the gate to change in
  response to changes in the input signals:
  \[t_{PD}=\left( t_{PHL} +t_{PLH}\right)/2 ,\] 
  where $t_{PHL}$ is the propagation delay time when the output changes from
  high to low, and $t_{PLH}$ that when the output changes from low to high.
\end{definition}
In general, given an $n$-input logic gate with delay $T$ and its inputs stable
at times $t_1,t_2,\ldots,t_{n}$, the earliest time in which the output is
stable is $max\left\{ t_1,t_2,\ldots,t_{n} \right\} +T$.

\begin{definition}[Circuit Delay]
  The \textbf{circuit delay} is the ``propagation of the entire circuit'', the
  longest time it takes for the input signals to propagate to the outputs.
\end{definition}
\begin{example}[Circuit Delay]
  Given a circuit, we can calculate the circuit delay by propagating the inputs
  forward past each logic gate or component to obtain the final delay.
  \begin{figure}[h!]
    \centering
    \includegraphics[width=0.5\textwidth]{circuit_delay}
  \end{figure}
\end{example}

\subsubsection{Alternative Logic Circuits}
\paragraph{Integrated Circuit}
\begin{definition}[Integrated Circuit]
  An \textbf{Integrated Circuit }(IC), also known as a chip or microchip, is a
  semiconductor device in which transistors, capacitors and resistors are
  fabricated.
\end{definition}
ICs can be categorised according to the number of components they contain,
\begin{itemize}
  \item \textit{Small Scale Integration}: Up to 100 electronic components
  \item \textit{Medium-Scale Integration}: From 100 to 3,000 electronic components
  \item \textit{Large-Scale Integration}: From 3,000 to 100,000 electronic components
  \item \textit{Very Large-Scale Integration}: From 100,000 to 1,000,000 electronic components
  \item \textit{Ultra Large-Scale Integration}: More than 1,000,000 electronic components
\end{itemize}
or according to its capabilities and limitations, into different logic families:
\begin{itemize}
  \item \textit{Transistor-Transistor Logic}: Uses bipolar junction transistors
  \item \textit{Complementary Metal-Oxide Semiconductor}: Uses metal-oxide
    semiconductor field effect transistors, consumes less power and has higher
    packing density, but is very sensitive to static electricity
  \item \textit{Emitter-Coupled Logic}: Uses bipolar circuit technology,
    designed for very high speed operations, but has high power consumption.
  \item among many others…
\end{itemize}

\begin{example}[TTL]
  The Transistor-Transistor Logic family is the largest, consisting various
  versions: \textit{standard TTL, low-power TTL, Schottky TTl, low-power
  Schottky TTL, advanced Schottky TTL,} etc. These ICs are encoded by four- to
  five-digit part numbers, with a two-digit prefix (usually 54 or 74) followed
  by a two- or three-digit code for the gate. Between the prefix and the gate
  code is a letter code indicating the TTL version.
\end{example}

\paragraph{Programming Logic Array}
\begin{definition}[Programming Logic Array]
  A \textbf{Programming Logic Array }(PLA) is a programmable integrated circuit
  that implements SOP circuits with multiple outputs using a 2-level AND-OR
  circuit.
\end{definition}
\begin{figure}[h!]
  \centering
  \includegraphics[width=0.3\textwidth]{PLA.png}
\end{figure}

\paragraph{Read Only Memory}
\begin{definition}[Read Only Memory]
  A \textbf{Read-Only Memory} is similar to PLA with a set of inputs
  (addresses) and a set of outputs, with a programmable mapping between inputs
  and outputs.
\end{definition}
A ROM is fully decoded and able to implement any mapping, in contrast to PLA
which may not have enough minterms.

\subsubsection{MSI Components}
In addition, we have various common medium-scale integration (MSI) components
that are used in circuit design, namely the \textit{decoders},
\textit{encoders}, \textit{demultiplexers}, and \textit{multiplexers}.

\paragraph{Decoders}
\begin{definition}[Decoders]
  A $n$-to-$m$-line, $n:m$, or $n\times m$ \textit{decoder}, converts binary
  information from $n$ input lines to $m\le 2^{n}$ output lines. Given an input
  of $k_{10}={\left( X_{n}X_{n-1}\cdots X_1 \right)}_2$, the decoder selects
  the $k$-th output line.
\end{definition}
\begin{remark}
  Each output is the minterm of all the variables, since every code uniquely
  identifies an output line to be selected.
\end{remark}
\begin{definition}[Enable]
  In these components, we may have an \textit{enable control} sigal, such that
  the device is only activated when the enable $E=1$. We say that the device is
  one-enable.
  \begin{figure}[h!]
    \centering
    \includegraphics[width=0.5\textwidth]{decoder_enable}
  \end{figure}\\
  Enables may be implemented with an AND gate for each of the outputs.
\end{definition}
\begin{definition}[Active High/Low Outputs]
  A device has normal outputs if it has active high outputs, and negated
  outputs if it has active low outputs.
\end{definition}
\begin{remark}
  Most Standard MSI Decoders are zero-enabled, and have active-low outputs.
  \begin{example}[74138 $3:8$ Decoder]
    We illustrate the representations of a standard 74138 decoder, a
    zero-enabled, active-low output, $3:8$ decoder.
    \begin{figure}[h!]
      \centering
      \includegraphics[width=0.5\textwidth]{decoder_74138}
      \includegraphics[width=0.4\textwidth]{decoder_74138_2}
    \end{figure}
  \end{example}
\end{remark}
Additionally, decoders can also be cascaded to construct larger decoders.
\begin{example}[$3:8$ Decoder]
  We can construct a $3:8$ decoder from two $2:4$ decoders with one-enable and
  an inverter.
  \begin{figure}[h!]
    \centering
    \includegraphics[width=0.3\textwidth]{decoder_cascade}
  \end{figure}
\end{example}
Decoders allow us to implement any function, by generating all the individual
minterms, and using an appropriate gate to implement the function. In general, for
a function $f=\sum m\left( a_1,a_2,\ldots,a_k \right)=\prod M\left(
b_1,b_2,\ldots,b_l \right)  $:
\begin{itemize}
  \item Decoder with active-high outputs
    \begin{itemize}
      \item with an OR gate:
        \[f=m_{a_1}+m_{a_2}+\cdots+m_{a_k}\]
      \item with a NOR gate:
        \[f=\left( m_{b_1}+m_{b_2}+\cdots+m_{b_l} \right)'\left[ =M_{b_1}\cdot
        \cdots \cdot M_{b_l} \right] \]
    \end{itemize}
  \item Decoder with active-low outputs
    \begin{itemize}
      \item with a NAND gate:
        \[f=\left( m_{a_1}'\cdot m_{a_2}'\cdot \cdots \cdot m_{a_k}\right)'\] 
      \item with an AND gate:
        \[f=m_{b_1}'\cdot m_{b_2}'\cdot \cdots \cdot m_{b_l}'.\] 
    \end{itemize}
\end{itemize}
They hence best lend themselves to circuits of many outputs, each a sum of a
few minterms. In general, any combinational circuit with $n$ inputs and $m$
outputs can be implemented with a $n:2^{n}$ decoder and $m$ logic gates. 

\paragraph{Encoders}
\begin{definition}[Encoders]
  A $m$-to-$n$-line, $m:n$, or $m\times n$ \textit{encoder}, converts binary
  information from $m<2^{n}$ input lines to $n$ output lines. If the $k$-th
  input line is active, the encoder outputs $\left( D_{n}D_{n-1}\cdots D_1
  \right)_2=k_{10}$. Behaviour is only defined for input lines of which exactly
  one is active.
\end{definition}
Encoders compute the inverse of decoders and can be implemented with $n$ OR
gates.
\begin{figure}[h!]
  \centering
  \includegraphics[width=0.4\textwidth]{encoder_truth}
  \includegraphics[width=0.4\textwidth]{encoder_circuit}
\end{figure}

\begin{definition}[Priority Encoders]
  A \textit{priority encoder} is an encoder with priority. If two or more
  outputs are active, the input with the highest priority takes precedence. If
  all inputs are inactive, behaviour is undefined.
\end{definition}

\paragraph{Demultiplexers}
\begin{definition}[Demultiplexers]
  Given an input line and a set of selection lines, a
  \textit{demultiplexer} directs data from the input to one selected output line
  by the control signals.
\end{definition}
\begin{figure}[h!]
  \centering
  \includegraphics[width=0.6\textwidth]{demultiplexer}
\end{figure}
The demultiplexer circuit is identical to a decoder with enable, by taking the
enable as the input and the input lines as the control.

\paragraph{Multiplexers}
\begin{definition}[Multiplexers]
  A \textit{multiplexer}, or data-selector, has $2^{n}$ input lines, $n$
  selection or control lines, and one output line. It directs data from the
  specified input line to the output line.
\end{definition}
It may be implemented with a $n:2^{n}$decoder, and getting each of the $2^{n}$
output's  product with the corresponding input line.

\begin{example}[74151A $8:1$ Multiplexer]
  We illustrate the representations of a standard 74151A multiplexer.
  \begin{figure}[h!]
    \centering
    \includegraphics[width=0.4\textwidth]{multiplexer}
    \includegraphics[width=0.4\textwidth]{multiplexer_2}
  \end{figure}
\end{example}

Multiplexers can also be cascaded to construct larger multiplexers.
\begin{example}[$8:1$ Multiplexer] We can construct a $8:1$ multiplexer from
  smaller multiplexers as shown.
  \begin{figure}[h!]
    \centering
    \includegraphics[width=0.4\textwidth]{multiplexer_cascade}
    \includegraphics[width=0.45\textwidth]{multiplexer_cascade_2}
  \end{figure}
\end{example}

Multiplexers also allow us to implement any function, by having each of the
variables as a selection line, and setting each input to be $1$ if it is the
$k$-th line, where $m(k)$ is a minterm of the function and $0$ otherwise. We
may also implement a $n$-variable function with a $2^{n-1}:1$ multiplexer, by
setting one of the variables to be input and the others the selection lines. We
do so by drawing the truth table, grouping by selection line values, and then
comparing the input variable to the function output.\\

We illustrate this with the following example, where $f(A,B,C)=\sum m\left(
0,1,3,6 \right) $.
\begin{figure}[h!]
  \centering
  \includegraphics[width=0.35\textwidth]{multiplexer_imp_2}
  \includegraphics[width=0.45\textwidth]{multiplexer_imp}
\end{figure}


\subsection{Combinational Circuits}
Within logic circuits, we first consider the class of combinational circuits,
where each output depends on only the immediate inputs. We may implement
combinational circuits with logic gates in the \textbf{gate-level design
method}, or functional blocks in the \textbf{block-level design method}.\\

\subsubsection{Gate-Level (SSI) Design}
In gate-level design, we follow the design procedure of (1) Stating the
problem, (2) Determining and labelling the inputs and outputs, (3) Drawing the
truth table, (4) Obtaining simplified Boolean expressions, (5) Drawing the
logic diagram.

\begin{example}[Full Adder] We demonstrate the design procedure with the Full
  Adder component.
  \begin{definition}[Half Adders]
    A \textit{Half Adder} is a circuit that adds 2 single bits $X$ and $Y$ to
    produce a result the carry-out bit $C$ and the sum bit $S$.
  \end{definition}
  From the truth table and algebraic simplification, we derive the simplified
  expressions for $C$ and $S$
  \[C=X\cdot Y,\qquad S=S'\cdot Y+X\cdot Y'=X\oplus Y.\]
  \begin{figure}[h!]
    \centering
    \includegraphics[height=0.15\textwidth]{half_adder_block}
    \includegraphics[height=0.15\textwidth]{half_adder_truth}
    \includegraphics[height=0.15\textwidth]{half_adder_circuit}
  \end{figure}\\
  We aim to obtain a full-adder to be able to add 3 single bits, which we call
  $X$, $Y$, and a carry-in bit $Z$. Using the K-map, we have the simplified SOP
  forms
  \[C=X\cdot Y+X\cdot Z+Y\cdot Z=X\cdot Y+\left( X\oplus Y \right)\cdot Z\]
  \[S=X'\cdot Y'\cdot Z+X'\cdot Y\cdot Z'+X\cdot Y'\cdot Z'+X\cdot Y\cdot Z=X\oplus\left( Y\oplus Z \right).\] 
  \begin{figure}[h!]
    \centering
    \includegraphics[height=0.2\textwidth]{full_adder_block}
    \includegraphics[height=0.2\textwidth]{full_adder_truth}
    \includegraphics[height=0.2\textwidth]{full_adder_circuit}
  \end{figure}\\
  We hence have the full adder constructed from two half-adders and an OR gate.
\end{example}

\subsubsection{Block-Level Design}
For more complex circuits, we may employ block-level design, where we decompose
the main problem to sub-problems recursively to be directly solved by blocks of
circuits.

\begin{example}[Parallel Adder] We demonstrate the block-level design procedure
  with the paralllel adder component.
  \begin{definition}[Parallel Adders]
    A \textit{parallel adder} is a circuit that adds two $4$-bit numbers $X=X_4
    X_3 X_2 X_1$ and $Y=Y_4 Y_3 Y_2 Y_1$ and a carry-in bit $C_1$ to produce a
    $5$-bit result $R=C_5S_4S_3S_2S_1$
  \end{definition}
  Rather than drawing a truth table with $9$-bit inputs, we notice the addition
  formula for each pair of bits with carry-in has the same function as a full
  adder
  \[C_{i+1} S_i=X_i+Y_i+C_i\]
  \[C_{i+1}=X_i\cdot Y_i+\left( X_i\oplus Y_i \right)\cdot  C_i,\qquad
  S_i=X_i\oplus Y_i\oplus C_i.\] 
  This allows us to cascade full adders via their carries to create a
  \textit{parallel adder} (or \textit{ripple-carry adder}).
  \begin{figure}[h!]
    \centering
    \includegraphics[height=0.2\textwidth]{parallel_adder_block}
    \includegraphics[height=0.2\textwidth]{parallel_adder_circuit}
  \end{figure}\\
\end{example}
Given the $4$-bit parallel adder, we may much more easily
implement a BCD to Excess-$k$ converter (by adding $k$ to a given BCD), or a
$4n$-bit parallel adder (by similarly cascading $n$ parallel adders).
\begin{remark}
  The circuit delay of a ripple-carry parallel adder is proportional to the
  number of bits it handles, since the carry bits has to propagate through each
  full adder.
\end{remark}

\begin{example}[Magnitude Comparator]
  Alternatively, to design a more complicated circuit, we may analyse the
  formulae or algorithms in its computation.
  \begin{definition}[Magnitude Comparators]
    A $n$-bit \textit{magnitude comparator} is a circuit that, given two
    $n$-bit numbers $A$ and $B$, compares and returns either $A>B$, $A=B$, or
    $A<B$.
  \end{definition}
  Given two unsigned values $A$ and $B$, we are only required to compare the
  MSB of the two values to conclude $A<B$ or vice versa, and check the next MSB
  only if the current MSB is equal. We hence may obtain the following circuit,
  where $A=A_3A_2A_1A_0$, $B=B_3B_2B_1B_0$, and $x_i=A_i\cdot B_i+A_i'\cdot
  B_i'$.
  \begin{figure}[h!]
    \centering
    \includegraphics[height=0.3\textwidth]{compa_block}
    \includegraphics[height=0.3\textwidth]{compa_circuit}
  \end{figure}
\end{example}


\subsection{Sequential Circuits}
Other than combinational circuits, we now consider the class of sequential
circuits, where each output depends not only on the present inputs, but also on
the current state.

\begin{definition}[Synchronous/Asynchronous Circuits]
  \textit{Synchronous circuits} are sequential circuits where outputs change
  only at a specific time, and \textit{asynchronous circuits} are that where
  outputs change at any time.
\end{definition}

\begin{definition}[Multivibrators]
  \textit{Multivibrators} are a family of sequential circuits, within which we
  distinguish bistable (2 stable states), monostable (1 stable state) and
  astable (no stable state) circuits.
\end{definition}

\begin{definition}[Memory Elements]
  A \textit{memory element} is a device which comprises state indefinitely, or
  can change state based on its inputs.
\end{definition}

The memory elements in a sequential circuit are implemented using bistable
logic devices, namely \textbf{latches} and \textbf{flip-flops}. Memory elements
can also be controlled by a \textbf{clock}, typically a square wave consisting
positive and negative edges, and positive pulses.

\subsubsection{Latches}
Latches are memory devices that are pulse-triggered. It is ON for the duration
of positive pulses, during which it would respond to inputs for state
transitions, and preserve its state when it is OFF

\paragraph{S-R Latch}
A S-R Latch takes in two inputs $S$ (SET) and $R$ (RESET) and provides two
complementary outputs $Q$ and $Q'$. The latch is in a SET state when $Q=$HIGH
and RESET state when $Q=$LOW.

In an active-high input S-R latch, the next state is determined by the input
carrying the high signal.
\begin{figure}[h!]
  \centering
  \includegraphics[width=0.4\textwidth]{sr_high}
  \includegraphics[width=0.35\textwidth]{sr_high_2}
\end{figure}\\
In an active-low input S-R latch, the next state is determined by the input
carrying the low signal.
\begin{figure}[h!]
  \centering
  \includegraphics[width=0.5\textwidth]{sr_low}
  \includegraphics[width=0.35\textwidth]{sr_low_2}
\end{figure}

A pulse-triggered latch gated S-R latch is a S-R latch with an enable input.
Outputs of the latch change only if the enable input is high.
\begin{figure}[h!]
  \centering
  \includegraphics[width=0.4\textwidth]{sr_gated}
\end{figure}
\vspace{-4mm}
\paragraph{D Latch}
The gated D latch eliminates the invalid command in the S-R latch by
constraining the inputs for an S-R latch to be $D$ and $D'$.
\begin{figure}[h!]
  \centering
  \includegraphics[width=0.4\textwidth]{d_gated}
  \includegraphics[width=0.3\textwidth]{d_gated_2}
\end{figure}
\vspace{-4mm}
\subsubsection{Flip-Flops}
Flip-flops are synchronous bistable memory devices that are edge-triggered.
That is, state is updated at a specified point on the clock signal, either at
the positive (rising) edge, or at the negative (falling) edge.
\begin{figure}[h!]
  \centering
  \includegraphics[width=0.65\textwidth]{flipflops}
\end{figure}
\vspace{-8mm}
\paragraph{S-R Flip-Flop} A positive-edge triggered S-R Flip-Flop behaves as a
synchronous version of an active-high input S-R latch. $Q^{+}=S+R'\cdot Q$ where $S\cdot R=0$.
\begin{figure}[h!]
  \centering
  \includegraphics[width=0.38\textwidth]{srff}
\end{figure}
\vspace{-8mm}
\paragraph{D Flip-Flop} A positive-edge triggered D flip-flop behaves as a
synchronous version of an active-high input D latch. $Q^{+}=D$.
\begin{figure}[h!]
  \centering
  \includegraphics[width=0.35\textwidth]{dff}
\end{figure}
\vspace{-8mm}
\paragraph{J-K Flip-Flop} In a J-K flip-flop, $Q$ and $Q'$ are fed back to the
pulse-steering NAND gates. We eliminate the invalid state, which becomes a toggle
state. $Q^{+}=J\cdot Q'+K'\cdot Q$.
\begin{figure}[h!]
  \centering
  \includegraphics[width=0.4\textwidth]{jkff}
  \includegraphics[width=0.38\textwidth]{jkff2}
\end{figure}
\vspace{-8mm}
\paragraph{T Flip-Flop} The T flip-flop is a single input version of the J-K
flip-flop where $J=K=T$, which triggers the toggle state. $Q^{+}=T\cdot Q'+T'\cdot Q$.
\begin{figure}[h!]
  \centering
  \includegraphics[width=0.35\textwidth]{tff}
\end{figure}
\vspace{-4mm}
\begin{definition}[Synchronous/Asynchronous Inputs]
  A \textit{synchronous input} is one where data is transferred to the output
  only at a specific time, and an \textit{asynchronous input} is one that
  affects the state independent of the clock.
\end{definition}
\begin{example}[Flip-Flops]
  For the flip-flops, we have the synchronous inputs $S$-$R$, $D$, and $J$-$K$.
  We may also have the asynchronous inputs $PRE$ or $SD$ (Preset/Direct Set),
  $CLR$ or $RD$ (Clear/Direct Reset). When $PRE=$HIGH, $Q=$HIGH
  \textit{immediately} and when $CLR=$HIGH, $Q=$LOW \textit{immediately}.
  \begin{figure}[h!]
    \centering
    \includegraphics[width=0.55\textwidth]{async_ff}
  \end{figure}
\end{example}

\subsubsection{Sequential Circuit Analysis}
Given a sequential circuit, we analyse its behaviour by deriving its state
table and state diagram. We do so with the following procedure:
\begin{enumerate}
  \item Derive flip-flop input functions for each flip-flop. We denote each
    flip-flop input with two letters, the input's name and the flip-flop's
    name. We algebraically describe the inputs to the flip-flops.
  \item Using the characteristic tables, obtain state equations for each
    flip-flop.
  \item Derive output functions for each non-flip-flop output of the circuit.
  \item Draw and fill the state table consisting all possible binary
    combinations of present states and inputs, to show the outputs and next
    state. This consists $2^{m+n}$ rows for $m$ flip-flops and $n$-inputs.
  \item The state table can be optionally drawn more compactly by representing
    the current state of multiple variables in a single column.
  \item Derive the state diagram, where each of the $2^{m}$ node denotes a
    state, an arrow denotes a transition of state and a label of the form $a/b$
    is attached to an arrow where $a$ denotes the inputs and $b$ the outputs of
    the circuit in that transition.
\end{enumerate}
\begin{example}[Sequential Circuit Analysis]
  \begin{figure}[h!]
    \centering
    \includegraphics[width=0.4\textwidth]{seq_an}
  \end{figure}
  With the given sequential circuit, we obtain the flip-flop input functions
  for flip-flops $A$ and $B$:
  \[DA=A\cdot x+B\cdot x,\qquad DB=A'\cdot x.\] 
  Using the characteristic table of the $D$ flip-flop, we hence have the state
  equations and output function of $y$
  \[A^{+}=A\cdot x+B\cdot x,\qquad B^{+}=A'\cdot x,\qquad y=\left( A+B \right)
  \cdot x'.\]
  With these equations and functions, we obtain the state tables and diagram of
  the circuit.
  \begin{figure}[h!]
    \centering
    \includegraphics[width=0.33\textwidth]{seq_an1}
    \includegraphics[width=0.3\textwidth]{seq_an2}
    \includegraphics[width=0.33\textwidth]{seq_an3}
  \end{figure}
\end{example}
\vspace{-8mm}
\subsubsection{Sequential Circuit Design}
\vspace{-4mm}
\begin{figure}[h!]
  \centering
  \includegraphics[width=0.5\textwidth]{seq_des}
\end{figure}
\vspace{-4mm}
Given a set of specifications in the form of state equations, state table, or
state diagram, we may derive a logic circuit fulfilling the specifications with
the following procedure:
\begin{enumerate}
  \item Derive the state table of the sequential circuit.
  \item Reduce states if necessary, removing duplicate states having
    identical outputs and next states.
  \item Assign a unique binary value to each state, if the given specifications
    depict states as labels (such as $X$, $Y$, $Z$)
  \item Determine the number of flip-flops required and label them. In general,
    a circuit with $n$ states require $\lceil{\log_2 n}\rceil$ flip-flops
  \item Using the excitation tables and the state tables, choose the
    appropriate type(s) of flip-flops to be used.
  \item Simplify the expressions of the flip-flop input functions and outputs.
  \item Draw the logic diagram of the circuit.
\end{enumerate}
\subsection{Memory}
Given these sequential circuits, we may implement a memory unit of $2^{k}$
words, each word consisting $n$ bits. The memory unit takes in $n$ data input
lines, a $k$-bit address and a $Read/\bar{Write}$ control signal, and provides
a $n$-bit data output.\\
\begin{figure}[h!]
  \centering
  \includegraphics[width=0.3\textwidth]{memory_unit}
\end{figure}
To conduct a Write operation, we have to transfer the address of the desired
word to the address lines, the data bits to be stored to the data input lines,
and the Write control line. Similarly, to conduct a Read operation, we have to
transfer the address of the desired word to the address lines, and activate the
Read control line.

\paragraph{RAM} In general, there are two types of RAM: static RAMs that uses
flip-flops as memory cells, and dynamic RAMs that use capacitor charges to
represent data. DRAMs are simpler in circuitry but have to be constantly
refreshed.
\begin{figure}[h!]
  \centering
  \includegraphics[width=0.49\textwidth]{ram}
  \includegraphics[width=0.5\textwidth]{ram2}
\end{figure}\\
A $4\times 3$ RAM (4 addresses and word size 3) can be constructed with a
circuit with a decoder and OR gates for outputs.
\begin{figure}[h!]
  \centering
  \includegraphics[width=0.4\textwidth]{ram3}
  \includegraphics[width=0.5\textwidth]{ram4}
\end{figure}\\
With smaller memory units, we can build bigger memory with an array of RAM
chips, by using decoders to select the correct RAM chip using the Chip select
control signals. This arrangement allows us to use more decoders and RAM chips
to arbitrarily expand our RAM storage.
\begin{figure}[h!]
  \centering
  \includegraphics[width=0.7\textwidth]{ram5}
\end{figure}
